\section{Experimental Method}
\label{sec:experiments}
\subsection{Populations and comparison conditions}
We combine retained component-study records from the CAMP artifact~\cite{campArtifact}
with new live Bowtie2 executions. Replotting the retained measurements does
not imply that their workflows were rerun. The feature and mmap studies use 33,516
tasks from GATK, Minimap2, Sarek, Seqinspector, Taxprofiler, and
Viralmetagenome. The fixed seed-1996 split assigns 28,490 tasks to
development and 5,026 to evaluation, using chronological ordering within
each workflow. Four feature views use static information and peak-memory
history (A+P), add consumption history ($\mathbf{h}_c$), add predicted
consumption ($\hat c$), or include both (full CAMP).

\begin{figure*}[!t]
\centering
\includegraphics[width=\textwidth]{\campReadyPath result_assets/audit_tradeoffs.pdf}
\caption{Two separate auditing studies. (a) Native downstream underallocation-rate difference over 17,067 tasks after the first update; negative values favor selective auditing. Bars are recorded paired 95\% bootstrap intervals. (b) Estimated net saving on 5,026 tasks with published allocations fixed (5--30\% budgets shown); bars are 95\% bootstrap intervals of five seed-level means. Allocation savings in (b) are not caused by the budget choices in (a).}
\label{fig:audit}
\end{figure*}

\label{sec:training-calibration}
The retained LightGBM ensembles and development-selected allocation
policies are preserved. Table~\ref{tab:features} compares each feature view
under the same global or process-tail configuration; it is a control for
feature interpretation, not a new routing-policy sweep. The original
endpoint comparison uses global A+P and process-tail CAMP. Published
requests and their retained reproduction are reported separately. The
mmap study uses paired task identities and peaks; each arm selects its
own allocation policy on development data, without test-time request
scaling. MdAPE denotes the median absolute percentage error of the median
peak prediction, not the error of the final reservation.

\subsection{Selective feedback and audit accounting}
The downstream replay uses a separate 30,478-task cohort: 9,141 development
and 21,337 evaluation tasks. Five fixed batches reveal recorded outcomes
at batch completion. Selective and uniform-random auditing use equal
per-workflow counts at budgets of 5, 10, 20, 30, and 50\%, with five
sampling seeds (1996--2000) and fixed model seeds. The primary evaluation
covers batches 2--5: 17,067 tasks after the first update. This is a batched
replay, not a reconstruction of live concurrency. Shared initial audited
history is excluded from incremental audit budgets.

A separate accounting study holds the published 5,026-task allocations
fixed and charges positive paired auditing overhead:
\begin{equation}
\begin{aligned}
 C_B&=\sum_{t:q_t=1}\hat y_t\max(d_t^{\rm audit}-d_t^{\rm base},0),\\
 G_B&=(W_{\rm AP}-W_{\rm CAMP})-C_B.
\end{aligned}
\label{eq:audit_cost}
\end{equation}
Memory and time are converted to GiB and hours. All tasks have paired
audit timings, but allocation-saving runtimes exist for only 4,577 tasks;
the 449 missing runtimes retain zero weight in the published saving.
$G_B$ excludes prediction, retraining, and retry costs and does not estimate
the allocation improvement caused by each budget. Recorded 95\% intervals
use paired grouped bootstrap for downstream differences and bootstrap of
five seed-level means for audit accounting.


\subsection{Live deployment and new-input branches}
\label{sec:live-method}
The Warm-Transfer study~\cite{campLiveArtifact}
has three stages:
cold, warm on the same inputs, and new input. Each contains 6,000 logical
tasks across seven Bowtie2 processes and produces 1,974 aligned chunks.
The original manifest selects 38 paired yeast ChIP-seq runs from ENA study
PRJEB11501; the new manifest selects 38 paired runs from PRJNA486148~\cite{enaOriginal,enaNew}.
The new runs were selected using input metadata and compressed size,
without execution outcomes: 15.20 versus 15.14~GiB. Their run/sample
identities, URLs, and FASTQ MD5 values do not overlap the original study.
The inputs are not distribution-matched: 34 new runs have 75-base reads
and four have 102-base reads, versus 102-base reads in the original set.

Cold execution starts from the fixed seed-1996 offline bundle trained on
28,490 non-Bowtie2 tasks and its leave-one-workflow-out cold-start policy.
It begins with no Bowtie2 history, but learns during execution. Warm starts
from the validated cold terminal model and matching history and calibration;
new input starts from the corresponding warm terminal state. Each stage
uses a fresh workflow directory, without cross-stage task-cache reuse.
The Direct-Transfer ablation instead starts the same new manifest from the
\emph{exact saved Warm-Transfer cold state}, skipping the original-data warm pass.
It does not rerun cold execution. Thus, the branches enter new input with
12,000 versus 6,000 prior Bowtie2 observations, respectively.

All live tasks are audited; selective auditing is evaluated only in the
separate replay. Successful completed outcomes become eligible for fitting;
failed outcomes remain diagnostics. The asynchronous learner first updates
after 128 eligible outcomes and then at 128-observation thresholds. The
latest 64 observations are held out from fitting for version-matched
calibration; a terminal update includes remaining eligible outcomes.
Consumption is measured after completion, whereas the task's allocation
uses a prelaunch prediction. Model seeds, allocation rules, and worker
settings are shared between branches.

Execution uses three 16-vCPU, 64-GiB workers (50~GiB exposed per worker to
Slurm). The main experiment uses a 16-vCPU, 64-GiB control node; the direct
cold branch uses an 8-vCPU, 32-GiB control node after downsizing. Both learners
retain four CPUs. Nextflow 25.10.7 launches Bowtie2 2.4.2 and samtools 1.20.
Slurm grants add a fixed 128~MiB to each CAMP request for execution overhead;
this allowance is excluded from $\hat y_t$, reported coverage, and $W$.
Task retries are limited to three after the initial attempt. Validation
checks logical-task coverage, audit outcomes, BAM integrity, and the
consistency of the terminal model, history, and calibration.

\subsection{Interpretation and validity}
Logical-task metrics pair the first CAMP request with the final successful
peak and runtime. Failed-attempt resource costs, auditing, prediction, and
retraining costs are excluded from $W$; actual grants and attempt records
are retained separately. A zero OOM count under the execution allowance
therefore does not establish that CAMP's nominal requests are sufficient.
The shared cold checkpoint avoids variation from an independent cold
rerun, but asynchronous completion order can change later model updates.
Control-node downsizing also prevents a controlled runtime comparison.
The model's system-context features retain the original deployment encoding
(13 nodes, 32 vCPUs and 244.14~GiB per node), rather than being refitted to
this cluster; conclusions are conditional on this configuration.
Only one execution of each branch and one new input study are evaluated,
so the results do not establish statistical significance or generalization
across studies.
